%% kapa: exam=1 round=28 year=2017 period=1 form=A subject=civil_law subject_ko=민법 seq=1 number=1
%% source: Q-Net 감정평가사 기출문제 (cst00309) artlSeq=5199603 file=extracted/1008/5199603/2017년 제28회 감정평가사 1차시험 문제지/1교시 A형.pdf
%% conversion: kordoc 4.24.51 -> tools/gen1.py (auto); review=proofread-sonnet; flags=none
\begin{question}{1}
법원(法源)에 관한 설명으로 옳지 않은 것은? (다툼이 있으면 판례에 따름)
\begin{choices}
  \item 사회생활규범이 관습법으로 승인되었다면 그것을 적용하여야 할 시점에서의 전체 법질서에 부합하지 않아도, 그 관습법은 법적 규범으로서의 효력이 인정된다.
  \item 법원은 관습법의 존부를 알 수 없는 경우를 제외하고 당사자의 주장ㆍ증명이 없어도 관습법을 직권으로 확정하여야 한다.
  \item 관습법은 법령과 같은 효력을 가지는 것으로서 법령에 저촉되지 않는 한 법칙으로서의 효력이 있다.
  \item 물권은 법률 또는 관습법에 의하는 외에는 임의로 창설하지 못한다.
  \item 강행규정 자체에 결함이 있거나 강행규정 자체가 관습에 따르도록 위임한 경우에는 사실인 관습에 법적 효력을 부여할 수 있다.
\end{choices}
\end{question}
\answer{1}
